\documentclass[10pt]{article}
\usepackage[utf8]{inputenc}
\usepackage[a4paper,height=20cm,width=13cm]{geometry}
\usepackage{amssymb}
\usepackage{dsfont}
\usepackage{calc}
\usepackage{graphicx}
% \usepackage{pst-node}
\usepackage{fourier}
\usepackage{euscript}
\usepackage{amsmath,amsthm}
\def\phi{\varphi}
\def\theta{\vartheta}
\def\P{\EuScript P}
\def\M{\EuScript M}
\def\B{\EuScript B}
\def\C{\EuScript C}
\def\D{\EuScript D}
\def\U{\EuScript U}
\def\V{\EuScript V}
\def\S{\EuScript S}
\def\sm{\smallsetminus}
\def\niff{\nleftrightarrow}
\def\ZZ{\mathds Z}
\def\NN{\mathds N}
\def\PP{\mathds P}
\def\QQ{\mathds Q}
\def\RR{\mathds R}
\def\<{\langle}
\def\>{\rangle}
\def\E{\exists}
\def\A{\forall}
\def\0{\varnothing}
\def\imp{\rightarrow}
\def\iff{\leftrightarrow}
\def\IMP{\Rightarrow}
\def\IFF{\Leftrightarrow}
\def\range{\textrm{im}}
\def\Mod{\textrm{Mod}}
\def\Aut{\textrm{Aut}}
\def\Th{\textrm{Th}}
\def\acl{\textrm{acl}}
\def\dom{{\rm dom}}
\def\eq{{\rm eq}}
\def\tp{\textrm{tp}}
\def\equivL{\stackrel{\smash{\scalebox{.5}{\rm L}}}{\equiv}}

\def\cnonfork{\mathbin{\raise1.8ex\rlap{\kern0.6ex\rule{0.6ex}{0.1ex}}\rlap{\kern1.1ex\rule{0.1ex}{1.9ex}}\raise-0.3ex\hbox{$\smile$} } }

\newcommand{\labella}[1]{{\sf\footnotesize #1}\hfill}
\renewenvironment{itemize}
  {\begin{list}{$\triangleright$}{%
   \setlength{\parskip}{0mm}
   \setlength{\topsep}{0mm}
   \setlength{\rightmargin}{0mm}
   \setlength{\listparindent}{0mm}
   \setlength{\itemindent}{0mm}
   \setlength{\labelwidth}{3ex}
   \setlength{\itemsep}{0mm}
   \setlength{\parsep}{0mm}
   \setlength{\partopsep}{0mm}
   \setlength{\labelsep}{1ex}
   \setlength{\leftmargin}{\labelwidth+\labelsep}
   \let\makelabel\labella}}{%
   \vspace*{-.3\baselineskip}
  \end{list}}
%\def\ssf#1{\textsf{\small #1}}
\newcounter{ex}
\newenvironment{exercise}{\bigskip\addtocounter{ex}{1}\textbf{Esercizio \theex.\ }}{}
\pagestyle{empty}
\parindent0ex
\parskip1ex
\raggedbottom
\def\nsR{{}^*\!\RR}
\def\ssf#1{\textsf{#1}}
\renewcommand{\baselinestretch}{1.3}


\usepackage{fancyhdr}
\pagestyle{fancy}
\lfoot{Teoria dei Modelli a.a.~2026/27}
\rhead{}
\rfoot{{\small\thepage}}
\cfoot{}

%
% INDICARE NOME E COGNOME DI TUTTI GLI AUTORI
%
\lhead{Nome/i Cognome/i}
%

\begin{document}

\begin{exercise}(Martina Apicella)
Prove that the following are equivalent
\begin{itemize}   
\item[1.] $T$ is $\omega$-categorical
\item[2.] $T$ is $\omega$-categorical over $A$ (for some/any) set $A$.
\end{itemize}
\end{exercise}

\begin{exercise}(Beatrice Calderaro)
Let $C$ be a finite set.
Prove that if $C\cap M\neq\varnothing$ for every model $M$ containing $A$, then $C\cap\acl A\neq\varnothing$.
\end{exercise}

\begin{exercise}(Francesco Dam)
Prove that the following are equivalent
\begin{itemize}   
\item[1.] $T$ is $\omega$-categorical
\item[2.] there is a countable model that is both saturated and atomic. 
\end{itemize}
\end{exercise}

\begin{exercise}(Filippo Giugie)
  Let $\phi(z)\in L(A)$ be a consistent formula.
  Prove that, if $a\in\acl(A,b)$ for every $b\models\phi(z)$, then $a\in\acl A$.
  Prove the same claim with a type $p(z)\subseteq L(A)$ for $\phi(z)$.
\end{exercise}

\begin{exercise}(Daniele Lanciotti)
  Let $M$ be $\omega$-saturated and $A\subseteq M$ finite.
  Suppose that every automorphism of $\U$ that fixes $M$ setwise fixes also $\acl A$ setwise.
  Prove that $\acl A\subseteq \acl\varnothing$.
  (The interesting case is when $M$ is countable.)
\end{exercise}

\begin{exercise}(Mattia Lo Mele \& Gabriele Maiolino)
  Let $\D\subseteq\U^x\times\U^y$. For any $a\in\U^x$ and $b\in\U^y$ we define $\D_{(a,y)}=\{y\,:\,(a,y)\in\D\}$ and $\D_{(x,b)}=\{x\,:\,(x,b)\in\D\}$.

  Assume $L(A)$ is countable.
  Let $\D\subseteq\U^x\times\U^y$ be such that for every $a\in\U^x$ and $b\in\U^y$ the sets $\D_{(a,y)}$ and $\D_{(x,b)}$ are definable over $A$.
  Prove that there is a model $M\supseteq A$ and a type $q(x,y)\subseteq L(A)$ such such that $q(M^x,M^y)=\D\cap \big(M^x\times M^y\big)$.

  Nota. Se non ci sono errori, questo dovrebbe essere un possibile analogo del teorema di Namioka. Forse non è esattamente giusto, o forse vale qualcose di più forte. Forse è connesso con la nozione di quasi-definibilità, ma non vedo come.
\end{exercise}

\begin{exercise}(Andrea Morello)
Let $a\in\U\sm\acl\varnothing$.
Prove that $\U$ is isomorphic to some $\V\preceq \U$ such that $a\notin\V$.
\end{exercise}
\end{document}


